\subsection{Module 10: RegA, PKA and ERK2 Brake}\label{sec:m10}

\paragraph{Role.}
The intracellular phosphodiesterase RegA (DdPDE2) \mk{M10-rega-pde}{degrades intracellular cAMP}~\cite{bader2007}.
It carries a phospho-relay receiver domain: \mk{M10-relay}{phosphorylation of Asp212 by the histidine phospho-transfer protein RdeA}~\cite{thomason1998,thomason1999} \mk{M10-rega-20}{activates the enzyme at least 20-fold}~\cite{thomason1999}, and \mk{M10-rega-km}{its cAMP-hydrolysing domain has a Michaelis constant of about 5\,$\mu$M}~\cite{thomason1998}.
\mk{M10-rega-loss}{Loss of RegA activates PKA and disturbs differentiation}~\cite{shaulsky1998}.
cAMP oscillations in populations are generated by a circuit in which \mk{M10-circuit}{cAMP activates PKA, PKA inhibits ERK2, and ERK2 controls cAMP levels through RegA}~\cite{laub1998,maeda2004}; \mk{M10-erk-phase}{the ERK2 activation oscillates in phase with the cAMP peaks}~\cite{laub1998,maeda2004}.
PKA also \mk{M10-pka-erk}{shapes ERK2 activation and adaptation}~\cite{aubry1997}.
Module~10 is the resulting brake: at rest RegA holds intracellular cAMP near zero, an ERK2 spike disarms it so that cAMP can accumulate, and its recovery and the PKA arm terminate the pulse and set the refractory tail.

\begin{table*}[!t]
\centering
\caption{Module~10 reactions: the intracellular brake. RegA is a conserved three-state moiety: RegA$_u$ (unphosphorylated, weakly active), RegA$_p$ (phosphorylated at D212, active) and RegA$_{pi}$ (phosphorylated by ERK2, inactive). PKA and ERK2 are the kinases of the Laub--Loomis/Maeda circuit.}
\label{tab:m10rxn}
\footnotesize
\renewcommand{\arraystretch}{1.15}
\begin{tabularx}{\textwidth}{@{}l>{\raggedright\arraybackslash}p{0.37\textwidth}>{\raggedright\arraybackslash}X@{}}
\toprule
ID & Reaction & Propensity (rate law)\\
\midrule
10.B1 & $\mathrm{RegA}_u \to \mathrm{RegA}_p$ & $k_{B1}[\mathrm{RegA}_u]$ \\
10.B2 & $\mathrm{RegA}_p \to \mathrm{RegA}_u$ & $k_{B2}[\mathrm{RegA}_p]$ \\
10.B3 & $\mathrm{RegA}_p+\mathrm{ERK2}^* \to \mathrm{RegA}_{pi}+\mathrm{ERK2}^*$ & $k_{B3}[\mathrm{RegA}_p][\mathrm{ERK2}^*]$ \\
10.B4 & $\mathrm{RegA}_{pi} \to \mathrm{RegA}_p$ & $k_{B4}[\mathrm{RegA}_{pi}]$ \\
10.C1 & $\mathrm{cAMP}_i+\mathrm{RegA}_p \to \mathrm{RegA}_p$ & $k_\mathrm{cat}[\mathrm{cAMP}_i][\mathrm{RegA}_p]$ \\
10.C2 & $\mathrm{cAMP}_i+\mathrm{RegA}_u \to \mathrm{RegA}_u$ & $\phi\,k_\mathrm{cat}[\mathrm{cAMP}_i][\mathrm{RegA}_u]$ \\
10.D1 & $\mathrm{cAMP}_i+\mathrm{PKA} \to \mathrm{cAMP}_i+\mathrm{PKA}^*$ & $k_{D1}[\mathrm{cAMP}_i][\mathrm{PKA}]$ \\
10.D2 & $\mathrm{PKA}^* \to \mathrm{PKA}$ & $k_{D2}[\mathrm{PKA}^*]$ \\
10.D3 & $\mathrm{PKA}^*+\mathrm{ERK2}^* \to \mathrm{PKA}^*+\mathrm{ERK2}$ & $k_{D3}[\mathrm{PKA}^*][\mathrm{ERK2}^*]$ \\
\bottomrule
\end{tabularx}
\end{table*}

\begin{table*}[!t]
\centering
\caption{Module~10 constants.}
\label{tab:m10const}
\footnotesize
\renewcommand{\arraystretch}{1.15}
\begin{tabularx}{\textwidth}{@{}l r l c >{\raggedright\arraybackslash}X@{}}
\toprule
Symbol & Value & Unit & Basis & Meaning, justification, source\\
\midrule
$[\mathrm{RegA}]_\mathrm{tot}$ & \num{1540} & $\mathrm{molec/voxel}$ & E & Copy number never measured; $0.32\,\mu$M puts RegA on the scale of ERK2 and PKA and keeps the required catalytic efficiency sub-diffusion-limited. \\
$[\mathrm{PKA}]_\mathrm{tot}$ & \num{1540} & $\mathrm{molec/voxel}$ & E & $0.32\,\mu$M. \\
$k_{B1}$ & \num{0.05} & $\mathrm{s^{-1}}$ & E & \mkt{M10-relay}{Lumped phospho-relay (RdeA$\to$RegA D212) that activates RegA}~\cite{thomason1999,thomason1998}. \\
$k_{B2}$ & \num{5.6e-3} & $\mathrm{s^{-1}}$ & C & \mkt{M10-relay}{Reverse phospho-transfer}~\cite{thomason1999}; chosen for a resting active fraction of 0.9. \\
$k_{B3}$ & \num{8} & $\mu\mathrm{M}^{-1}\mathrm{s}^{-1}$ & C & \mkt{M10-erk-rega}{ERK2$^*$ phosphorylates and thereby inhibits RegA}~\cite{laub1998,maeda2004}. The Laub--Loomis coefficient (\mkt{M10-laub-k8}{$k_8=1.3$ in relative units}~\cite{laub1998}) is not a molecular rate; the value gives $\tau\approx0.5$\,s at the ERK2$^*$ peak so that the brake is released inside the X spike. \\
$k_{B4}$ & \num{0.0167} & $\mathrm{s^{-1}}$ & C & The reset reaction ($\tau=60$\,s) that ends each pulse; of the order of the RegA replenishment coefficient \mkt{M10-laub-k7}{$k_7=2\,\mathrm{min^{-1}}$ of the Laub--Loomis model}~\cite{laub1998} and set so that the cell has reset within about 6\,min. \\
$k_\mathrm{cat}$ & \num{30} & $\mu\mathrm{M}^{-1}\mathrm{s}^{-1}$ & D & $k_\mathrm{cat}/K_M=3\times10^{7}\,\mathrm{M^{-1}s^{-1}}$: $k_\mathrm{cat}\approx150\,\mathrm{s^{-1}}$ at \mkt{M10-rega-km}{the measured $K_M\approx5\,\mu$M}~\cite{thomason1998}. \\
$\phi$ & \num{0.05} &  & L & Activity of RegA$_u$ relative to RegA$_p$; \mkt{M10-rega-20}{phosphorylation activates RegA by at least 20-fold}~\cite{thomason1999}. \\
$k_{D1}$ & \num{0.25} & $\mu\mathrm{M}^{-1}\mathrm{s}^{-1}$ & E & cAMP binding to the PKA regulatory subunit; with $k_{D2}$ this gives $K_d=k_{D2}/k_{D1}=0.1\,\mu$M (model choice; \mkt{M10-circuit}{PKA activation is the second element of the circuit of}~\cite{laub1998}). \\
$k_{D2}$ & \num{0.025} & $\mathrm{s^{-1}}$ & D & PKA$^*$ decay ($\tau=40$\,s), from \mkt{M10-laub-k4}{the Laub--Loomis coefficient $k_4=1.5\,\mathrm{min^{-1}}$}~\cite{laub1998}; sets the minutes-long refractory tail. \\
$k_{D3}$ & \num{0.0417} & $\mu\mathrm{M}^{-1}\mathrm{s}^{-1}$ & D & PKA$^*$ inhibits ERK2, from \mkt{M10-laub-k6}{the coefficient $k_6=0.8$ of}~\cite{laub1998} (relative units, normalised here by the PKA pool); PKA also \mkt{M10-pka-erk}{shapes ERK2 activation and adaptation}~\cite{aubry1997}. \\
$D_{\mathrm{RegA}},D_{\mathrm{PKA}},D_{\mathrm{ERK2}}$ & \num{10} & $\mu\mathrm{m^2/s}$ & E & Cytosolic value. \\
\bottomrule
\end{tabularx}
\end{table*}


\paragraph{Reactions.}
\begin{description}\itemsep0pt
\item[10.B1/10.B2] The phospho-relay activates RegA and reverses. The pair sets the resting active fraction, 0.9. Only \mk{M10-relay}{the lumped transfer from RdeA to RegA} is modelled~\cite{thomason1999}.
\item[10.B3] ERK2$^*$ converts active RegA into an inactive, ERK2-modified form. The model implements \mk{M10-erk-rega}{the control of RegA by ERK2}~\cite{maeda2004} as inactivation, the sign required for an ERK2 spike to escape the brake.
\item[10.B4] The inactive form returns to the active form ($\tau=60$\,s), since \mk{M10-rega-phosphatase}{a protein phosphatase activates RegA}~\cite{laub1998}. This reset step is the one that terminates each pulse; without it RegA would be lost irreversibly.
\item[10.C1/10.C2] RegA hydrolyses cAMP$_i$. The rate is linear in cAMP because cAMP$_i$ stays far below the Michaelis constant; \mk{M10-rega-20}{the unphosphorylated form keeps $1/20$ of the activity of the phosphorylated form}~\cite{thomason1999}.
\item[10.D1/10.D2] \mk{M10-pka-r}{cAMP activates PKA by binding to its regulatory subunit}~\cite{laub1998}, and PKA$^*$ decays with $\tau=40$\,s, which sets the minutes-long tail of the refractory period.
\item[10.D3] \mk{M10-circuit}{PKA$^*$ inactivates ERK2$^*$}, which closes the loop~\cite{laub1998,aubry1997}.
\end{description}

\paragraph{Constants.}
The Laub--Loomis coefficients are rates of a normalised-activity model and not molecular rate constants, so they cannot be used literally.
The Laub--Loomis coefficients refer to activities in relative units, so they cannot be used as molecular rate constants.
In the calibration of this model, a rate for RegA inactivation taken over literally kept the brake closed (knock-down $1.3\times$); the present value gives a knock-down of RegA$_p$ of $64\times$.
The cyclase-RegA competition needs $k_\mathrm{cat}/K_M\approx3\times10^{7}\,\mathrm{M^{-1}s^{-1}}$ with the RegA pool of $0.32\,\mu$M, which is sub-diffusion-limited; an earlier pool of $0.032\,\mu$M would have needed a catalytic efficiency near the diffusion limit before RegA could compete with export.
The copy numbers of RegA, PKA and ERK2 have not been measured.
Only $k_{D2}$ and $k_{D3}$ are tied to coefficients of the Laub--Loomis model ($k_4=1.5$ and $k_6=0.8$ in relative units, the latter divided by the PKA pool); $k_{B4}$ is calibrated, and the resulting pulse (an export of about $10^7$ molecules, a RegA knock-down of $64\times$ and a full reset within about 6\,min) is what the module is calibrated to.
